% 副露（フーロ）：他のプレイヤーの打牌を取得することにより面子を完成させること。「鳴く」ともいう

\documentclass[../nankiru_main.tex]{subfiles}

% override main tex render option
\renewcommand{\renderOption}{1}

\begin{document}

\subfilepreamble{副露判断}

下記の条件がどちらの一つが満たす手牌は鳴くべき\citep[][pp. 36--41, p. 45, pp. 58--61]{hirasawa-naki}、

\begin{itemize}
\item 立直の打点上昇が限られた
\item 愚形（または枯れている両面）を解消するため
\item 手牌の価値が満貫ある
\end{itemize}

下記の条件を全部満たす手牌は鳴きべきではない\multicitep{hirasawa-naki, p.~48; rb1, p.~216}、

\begin{itemize}
\item 安い（鳴いたら2000に過ぎない）
\item 遠い（頭がない・3副露の可能性が高い）
\item 安全牌がない
\end{itemize}

鳴きのメリットとデメリットは表~\ref{tab:naki-compare} の通り。同じ手牌でも場況によって鳴くべきかないかが分かれる。たとえば、アガリトップ、南場のトップ目、東風戦の場合、誰より早く和了るべき場面なら、後付けのリスクがあっても積極的に鳴く\citep[][pp. 138--145]{tetsusenryaku}。

\begin{table}[ht]
  \centering
  \begin{tabular}{lll}
    & 鳴き & 門前 \\
    \hline
    速度 & $\bigcirc$ & $\times$ \\
    打点 & $\times$ & $\bigcirc$ \\
    守備 & $\times$ & $\bigcirc$
  \end{tabular}
  \caption{鳴きのメリットとデメリット\citep[][p.~87]{hirasawa-naki}}
  \label{tab:naki-compare}
\end{table}

\subsubsection*{打点}

\textbf{両面＋両面の受けで鳴きにより打点が1/3になると中盤まで立直を狙う}\citep[][p.~33]{hirasawa-naki}、\textbf{鳴いても打点がそんなに下がらないなら中盤に入ると鳴く}\citep[][p.~38]{hirasawa-naki}。

\tehai{405667m 22667p 34s}[東1局　西家　ドラ \mj{3z}][naki-1]

例え\ref{mj:naki-1}の場合、門前で行けばピンフがなくても少なくとも5200点がある。一方、\mj{2568p25s} で鳴いたらこの手牌の価値は2000点に過ぎない。よって、12巡目までは鳴かない（もちろん相手の速度や待ちの残り枚数で調整することもよくある）。

\tehai{405667m 22667p 34s}[東1局　西家　ドラ \mj{4m}][naki-2]

手牌\ref{mj:naki-2}にもう一つのドラがあって、鳴いたら最低限の価値が3900点になる。中盤まで立直を狙えばその打点の期待値は跳満になければならない。よって、和了率を優先するなら6巡目に入ったら鳴く\footnote{こちらは一向聴の場合。両向聴ならもっと早く鳴けばいいと思う。}。

\tehai{3406m 12456p 3367s}[東1局　西家　ドラ \mj{3s}][naki-3]

\textbf{満貫があるなら全部鳴く}\citep[][p.~61]{hirasawa-naki}。手牌\ref{mj:naki-6}はタンヤオ・ドラ3でほぼ満貫。\mj{12p} の辺張を生かす役は立直しかない。タンヤオを確定ため向聴数を進めなくてもチー \mj{2456m58s} あと \mj{12p} を落とすかチー \mj{3p}（\mj{3'24p} の形で）にするか、満貫に全速前進。

\subsubsection*{愚形解消}

\textbf{ネックである愚形が解消されたら全てを鳴く}\citep[][p.~40]{hirasawa-naki}。愚形＋両面なら両面を中盤後半までスルーする方がいい。例え手牌\ref{mj:naki-3}が愚形＋両面の受けがあって、\mj{4s} が何巡目で来てもチーテンを取る。\mj{58p} なら9巡目までスルー。

\tehai{405667m 2267p 35s 2z}[東1局　北家　ドラ \mj{3z}][naki-4]

\textbf{ネックではない愚形を鳴かない}\citep[][p.~40]{hirasawa-naki}。特にブロックが６つあって、愚形が一つだけで、三色などの役を確定するためではないと鳴かない。よって、\ref{mj:naki-5}は \mj{7m} を鳴かないが、\ref{mj:naki-6}は打点もう十分なので \mj{7m} はチーする。

\tehai{3468m 23488p 3467s}[東1局　西家　ドラ \mj{2z}][naki-5]

\tehai{2468m 23488p 3468s}[東1局　西家　ドラ \mj{8p}][naki-6]

手牌\ref{mj:naki-20}のように愚形ばかりで向聴数を進めないけど \mj{2m} をポンすることにより受け入れが \mj{6m3p} から \mj{567m234p} へ変化するのでここはポンすべき\citep[][p.~199]{hirasawa-naki}。手牌\ref{mj:naki-21}のように片方が愚形の場合でも頭をポン有利。手牌\ref{mj:naki-22}のように暗刻があるなら頭を積極的に消していくのが有利。

\tehai{12257m 24p 789s -5'55z}[東1局　西家　ドラ \mj{5z}][naki-20]

\tehai{12257m 45p 789s -5'55z}[東1局　西家　ドラ \mj{5z}][naki-21]

\tehai{12257m 45p 888s -5'55z}[東1局　西家　ドラ \mj{5z}][naki-22]

\subsubsection*{聴牌までの距離} 

\textbf{鳴くと頭の候補が2枚減るので頭がないと鳴きを控えよう}\citep[][p.~55]{hirasawa-naki}。但し、\mj{2345m} や \mj{23445m} のような頭が出来やすい形があれば頭がなくても大丈夫。

\tehai{3468m 3468p 067s 34z}[東1局　西家　ドラ \mj{2z}][naki-7]

\tehai{3468m 35p 34056s 34z}[東1局　西家　ドラ \mj{5s}][naki-8]

上家が \mj{7m} を切った。手牌 \ref{mj:naki-7} で \mj{7m} を鳴いたら頭の候補が \mj{68m} 2種減る。手牌 \ref{mj:naki-8} で鳴いてもツモ \mj{2356s} により頭ができる上、最後まで頭が出来なくても \mj{25s} 待ちや \mj{36s} 待ちにすることもできる。

\subsubsection*{相手の速度} 

\textbf{相手が聴牌している（または聴牌していそう）なら副露して和了率を優先する}\citep[][p.~153]{hirasawa-naki}。上記のような中盤の終わりまでスルーの手牌で、相手がもう立直をかけている、もしくはもう両副露しているなら、手牌を鳴いて早く聴牌を取るのは得。

\subsubsection*{上家の捨て牌の傾向} 

\textbf{上家が捨てそうな牌なら鳴く}\citep[][p.~158]{hirasawa-naki}。特に上家が混一色を狙って、他の色の牌を切らなければならない状況。

\tehai{11m 130689p 1346s 5z}[東1局　北家　6巡目　ドラ \mj{3m}][naki-9a]

手牌\ref{mj:naki-9a}はバラバラの手牌。6巡目でマンズの混一色をしている上家が \mj{0s} を切った。この手を上がる役はピンズのイッツーだけがチー \mj{0s} により手牌の価値が高まる上上家がピンズを切ってくれる可能性が高いのでここでチー。

\tehai{68m 340688 4088s 3z}[東1局　北家　5巡目　ドラ \mj{3m}][naki-9b]

手牌\ref{mj:naki-9b}もバラバラの手牌。上家が \mj{5p} を切った。カン \mj{7m} を上家から鳴ける可能性が非常に低いので、ここで \mj{5p} をチーして \mj{68m} の嵌張を先に処理する。

%--------------------------------------------
\subsection{役牌}

\textbf{役牌があればほとんどの場合は鳴く}。鳴かないの場合は次の二つ\citep[][p.~103]{hirasawa-naki}、

\begin{itemize}
\item 頭が役牌しかない
\item 鳴いたところで安くて遠い
\end{itemize}

\subsubsection*{頭が役牌しかない} 

\tehai{3457m 67p 2346s 266z}[東1局　西家　ドラ \mj{3m}][naki-7a]

例え\ref{mj:naki-7a}のように頭が \mj{6z} の対子しかない場合、\mj{6z} 鳴いても速度が上がってないのでスルー。2枚目が切られてもスルー\citep[][p.~104]{hirasawa-naki}。

\tehai{345m 667p 23467s 66z}[東1局　西家　ドラ \mj{3m}][naki-7b]

逆に\ref{mj:naki-7b}のようにポンテンができる手牌では、立直・ドラ1の打点は最低2600点、発・ドラ1の2000点からあまり高くならないので中盤に入るとポンテンをとるの方がいい\citep[][pp.~105--106]{hirasawa-naki}。

\tehai{123m 23467p 89s 277z}[東1局　西家　5巡目　ドラ \mj{4p}][naki-7c]

手牌\ref{mj:naki-7c}のように残りの搭子が愚形なら、役牌対子を鳴いても愚形搭子の縦引きで受け入れが大きく下がらないなら鳴く。

\tehai{307m 4789p 207s 277z}[東1局　西家　4巡目　ドラ \mj{8p}][2-6-7]

手牌\ref{mj:2-6-7}は最低でも7700点。\mj{7z} をポンしたら頭がなくなるがまず役を確定する\citep[][pp.~164--167]{tetsusenryaku}。

\subsubsection*{鳴いたところで安くて遠い} 

\tehai{135m 4789p 268s 277z}[東1局　西家　ドラ \mj{7m}　3巡目][naki-7d]

三向聴・愚形ばかり・頭がないという遠い手牌で安くなら鳴かない\citep[][p.~110]{hirasawa-naki}。例として\ref{mj:naki-7d}はスルー。逆に価値があればほぼ全部鳴く。


%--------------------------------------------
\subsection{トイトイ}

\tehai{1167m 2244p 788s 77z}[東1局　西家　4巡目　ドラ \mj{8p}][2-8-1]

七対子一向聴で役牌対子があると鳴きに寄せる。対子の種類が12・89の牌が多いならなおさら。手牌\ref{mj:2-8-1}に \mj{44p} のネックがあって出たら積極的に鳴こう\citep[][pp.~204--211]{tetsusenryaku}。

\tehai{3m 3399p 11889s 224z}[東1局　西家　5巡目　ドラ \mj{3m}][2-8-2]

手牌\ref{mj:2-8-2}をトイトイにするとドラ \mj{3m} が働きないので手牌の価値は 2600 に過ぎない。親番で早く役を確定したい場面などではないなら鳴かない。

%--------------------------------------------
\subsection{片アガリ}

\textbf{片アガリの場合でも通常の副露判断に従う}\citep[][p.~123]{hirasawa-naki}。手牌の価値があれば振聴になるリスクを負って早く聴牌を取るの方がいい。

\tehai{111789m 5578p 78s 2z}[東1局　西家　ドラ \mj{1m}　6巡目][naki-10]

手牌\ref{mj:naki-10}の場合、上家が \mj{9p} を切った。両面＋両面ですが聴牌までにおよそ8巡目がかかる。\mj{9s} だけでロンできても早く満貫を確定の方がいい。

\tehai{136m 23p 120589s 25z}[東1局　西家　ドラ \mj{3z}　2巡目][naki-11]

\tehai{136m 23p 230589s 25z}[東1局　西家　ドラ \mj{3z}　2巡目][naki-12]

手牌\ref{mj:naki-11}と\ref{mj:naki-12}はバラバラの牌姿。聴牌までまた遠いの階段、両副露になる可能性が高い上、リーチが来て2000点の愚形で押し返すのは難しい。\textbf{「和了より押し引きの精度が要求されるルール」であればスルー}\citep[][p.~124]{hirasawa-naki}。しかも\ref{mj:naki-12}が両面二つあるので門前で聴牌する確率がより高い。\textbf{アガリトップなら鳴く}\citep[][p.~123]{hirasawa-naki}。

%--------------------------------------------
\subsection{タンヤオへの移行}

\textbf{3900点があるなら向聴数が変わらなくても速度が上がるため鳴く}\citep[][pp.~127--128]{hirasawa-naki}。2000点なら場況（相手の速度・待ちの状況・ドラが使えるかなど）で判断する。

\tehai{367m 234499p 3377s}[東1局　西家　ドラ \mj{3z}　5巡目][naki-13]

\tehai{367m 234499p 3377s}[東1局　西家　ドラ \mj{3s}　5巡目][naki-14]

\tehai{367m 234499p 3377s}[東1局　西家　ドラ \mj{7z}　5巡目][naki-15]

\tehai{367m 234499p 3377s}[東1局　西家　ドラ \mj{3m}　5巡目][naki-16]

手牌\ref{mj:naki-13}で相手が切った \mj{7s} をポンしたら僅かなタンヤオのみので、スルーの方がいい。逆に\ref{mj:naki-14}は3900点あるので鳴く。手牌\ref{mj:naki-15}は鳴いたら2000点だけで、アガリトップでなければ立直に寄る。手牌\ref{mj:naki-16}も2000点ですが、門前で行けば \mj{3m} が不要になる可能性が高い\footnote{特に \mjfn{9p37s} が暗刻になる場合}のでタンヤオ・ドラ1の方がいい。

%--------------------------------------------
\subsection{一発消しと海底ずらし}

\textbf{一発消しの有効局面はオーラス跳満ツモ条件だけ}\citep[][p.~188]{hirasawa-naki}。つまり、自分がリードしている、点差が14000点の二位の相手が立直をかけて来て、安全牌が足りているなら積極的に一発を消すべき。

\begin{rittai}[h]
  \centering
  
  \drawrittai{
    jicha/seat = 北,
    jicha/hand = 12206m 23556p 247s -8s,
    jicha/river = XXXXXX,
    jicha/riverb = XXXXXX,
    jicha/riverc = XXX,
    %
    toimen/river = 3z 9s 2m 2z 6z 4'm,
    toimen/riverb = 4^p 5^z 7^3^s 3^m 2^z,
    toimen/riverc = 2^s 7^z 1^2^m,
    %
    kami/river = XXXXXX,
    kami/riverb = XXXXXX,
    kami/riverc = XXXX,
    %
    shimo/river = XXXXXX,
    shimo/riverb = XXXXXX,
    shimo/riverc = XXXX,
    %
    dora = 2p
  }
  
  \caption{海底ずらしの場合}
  \label{rittai:haitei}
\end{rittai}

\textbf{海底ずらしは常にすべき}\citep[][p.~190]{hirasawa-naki}。特に自分の手に和了なさそうな場面。ただし局収支の差は僅かな100点\citep[][p. 188]{miinin-choice}。立体~\ref{rittai:haitei} では現物が足りているが、\mj{1m} を切ると \mj{3m} がチーできなくなって、\mj{27s} 切ると同じ上家が切った現物をチーできないので、ここは4枚見えていてポンできない \mj{2m} を切る。

%--------------------------------------------
\subsection{赤ドラを晒すと晒さない場面}

\tehai{4056m 2245p 3s 3z -5'55z}[東1局　西家　ドラ \mj{5m}　5巡目][naki-17]

\textbf{通常は赤ドラを晒すべきだが、黒5を晒すのは相手に警戒されたくない場面}\citep[][p.~194]{hirasawa-naki}。例え\ref{mj:naki-17}の副露手、\mj{5m} を晒すと相手から見える情報は「白ドラ1」ので、あまり警戒されない晒す方。

\tehai{4056m 246p 1136s 55z}[東1局　西家　ドラ \mj{5m}　5巡目][naki-18]

逆に\ref{mj:naki-18}のようなバラバラの手で、どうせ仕掛けが確定しているなら赤ドラを晒して相手の手を妨げるのが有効\footnote{ブラフは読みができて、または相手の副露に気をついている相手に有効。相手は読みができない、または何があっても気にせず全ツッパしているならブラフは効かない。}\citep[][pp.~196--197]{hirasawa-naki}。

\end{document}